%% ============================================================================
\section{The Framework in Detail}
\label{app:framework}
%% ============================================================================

This section expands \S\ref{sec:framework} step by step. Nothing in it is needed to check the
results, whose proofs are in App.~\ref{app:proofs}.

\subsection{The Committed World (\S\ref{sec:setup})}

Verdict polarity is our default for two reasons. It depends on the claims and the evidence, not on
the response, so all responses over the same claims are judged against the same world. And it does
not depend on how a response presents a claim, which a decomposition may not record. The verdict is
binary: a pipeline that separates refuted from unverifiable claims must map both to $0$, so a claim
the evidence fails to support is treated as false, and a refutation of it is credited as a
refutation of a falsehood. Assertion polarity needs the decomposition to record stance.

Verdict polarity makes the score a measure of how well an argument holds up \emph{given} the
verdicts, and four consequences follow.
\begin{itemize}[leftmargin=1.2em,itemsep=1pt,topsep=2pt]
\item A response that sets a true claim against a false one is credited for refuting the
  falsehood, even if it asserts both. Asserting the falsehood is penalized by the factuality score,
  not by this one; assertion polarity is the reading that treats such a pair as a
  self-contradiction.
\item An entailment from a false claim cannot be broken in the committed world, so reasoning from
  or among falsehoods is not penalized. It is even credited as a support, because it rules out a
  world other than $\xs$.
\item Relations are read as material constraints, so the score does not ask whether a premise is
  relevant to its conclusion (the miners emit only relations the response itself draws), and the
  soundness of premises is again the factuality score's concern.
\item The score inherits the verdicts' errors. A wrong verdict on a claim can flip the sign of the
  relations at that claim: always for equivalences and exclusions, and for the other couplings
  depending on the value of the other endpoint.
\end{itemize}
The score is therefore not the truth-independent coherence of a set of propositions studied in
Bayesian epistemology \citep{shogenji1999,olsson2002}. It asks how well the response's relations
agree with a given assignment of truth values, and assertion polarity is its closest analogue to
the truth-independent notion.

\subsection{Couplings, Senses and Signs (\S\ref{sec:taxonomy})}

\paragraph{Couplings and natural logic.} The couplings are the constraints that the basic semantic
relations of natural logic \citep{maccartney2009} impose on truth values: \ent{} is forward
entailment (reverse entailment is \ent{} with the scope reversed), \eqv{} equivalence, \con{}
alternation, \cnc{} cover, \exc{} negation, and \emph{none} independence. A coupling only rules
cells out. Unlike the natural-logic relations, it does not also require the remaining cells to be
possible, so \con{}, for instance, does not exclude that the two claims are also exhaustive. The set
of couplings is closed under negating an endpoint: negating the target of \ent{} gives \con{},
negating its source gives \cnc{}, negating both reverses it, and negating one end of \eqv{} gives
\exc{}.

\paragraph{The compile map.} It makes three deliberate commitments.
\begin{itemize}[leftmargin=1.2em,itemsep=1pt,topsep=2pt]
\item Five senses collapse onto \ent{}, because a pairwise potential over truth values cannot tell
  causation from evidence. The sense is kept for interpretation only, and every entailment is
  oriented from the supporting claim to the supported one, which for both causal senses means from
  cause to effect.
\item Only \sense{Refutation}, where the response asserts that two claims cannot both hold (``it
  cannot have, since\dots'', ``cannot be squared with''), compiles to \con{}. \sense{Contrast} and
  \sense{Concession} do not: in PDTB and RST both of their arguments are typically true (``sales
  rose, but profits fell''; ``although the airline blamed pilot error, the tribunal found a
  defect''), and what they deny is an expectation, not the conjunction.
\item Temporal senses produce no factor, since ordering two events says nothing about either
  one's truth, so reordering such claims cannot change the score.
\end{itemize}

\paragraph{Signs.} $\Rsup$ collects the relations that the committed world upholds and $\Rcon$
those it breaks. The split is a property of the relation graph \emph{together with} $\xs$: the
relation alone says which joint values it rules out, and the verdicts say whether the response's
own world is among them. It is a partition, since every relation that produces a factor is exactly
one of the two. The verdicts decide only a relation's sign; how much it earns depends on the rest
of the model (Lemma~\ref{lem:edit}).

\paragraph{The Voyager pair at the asserted world.} Both responses deny exactly the claims the
verdicts reject, so assertion polarity gives the same split as verdict polarity. The split would
change if the denied claims were read as asserted, at $\xs=\one$: every contradiction would become a
conflict, including A's two refutations, and every entailment a support, including B's
$a_1\Rightarrow a_4$ (App.~\ref{app:voyager}).

\paragraph{Claims and their negations.} A decomposition may record a denied claim as itself, as in
Figure~\ref{fig:voyager}, or as its asserted negation, ``the instrument had not failed''. The choice
does not matter (a posterior of exactly $\frac12$, which thresholding sends to $1$ either way,
aside). Replacing $a_i$ by $\neg a_i$ flips $\xs_i$, replaces $\pi_i$ by $1-\pi_i$, and rewrites
every coupling at $a_i$ by the closure above, so $a_2\con a_4$ becomes $a_2\Rightarrow\neg a_4$.
Each relation keeps its sign, and the score is unchanged (Prop.~\ref{prop:props}(iv)). What matters
is only that the verdict or stance and the coupling refer to the same form of each claim.

\subsection{Estimating Relations (\S\ref{sec:mining})}

The product of Eq.~\ref{eq:decomp} keeps a confidently identified weak link apart from a doubtfully
identified strong one. It is a modelling choice. A relation whose type is uncertain could instead
enter as a mixture of tables, but when the mixed couplings disagree in sign at $\xs$, the mixture
can place the committed cell strictly between its extremes (an entailment--contradiction mixture
whose support is the likelier reading does). That loses the guarantee of
Theorem~\ref{thm:necessity}, whereas the product keeps every factor two-valued.
App.~\ref{app:estimation} details both miners.

\subsection{The Coherence MRF (\S\ref{sec:mrf})}

% The coherence MRF of Response A (Fig. fig:voyager), as a factor graph: circles are
% the claim variables, grey squares the unary prior factors, black squares the
% sign-definite relation factors of Eq. eq:mrf. Numbers: numbers.tex.
\begin{figure}[t]
\centering
\begin{tikzpicture}[x=13mm, y=11mm]
  \node[atomtrue]  (a1) at (0,0)       {$a_1$};
  \node[atomtrue]  (a2) at (1.2,0)     {$a_2$};
  \node[atomtrue]  (a3) at (2.4,0)     {$a_3$};
  \node[atomfalse] (a4) at (1.2,-1.4)  {$a_4$};
  \node[atomfalse] (a5) at (3.0,-1.4)  {$a_5$};
  \foreach \v/\x/\y/\lab in {a1/0/0.85/$\phi_1$, a2/1.2/0.85/$\phi_2$, a3/2.4/0.85/$\phi_3$,
                              a4/1.2/-2.25/$\phi_4$, a5/3.0/-2.25/$\phi_5$}{
    \node[unary] (u\v) at (\x,\y) {}; \draw[flink] (u\v) -- (\v);
    \node[prlab,left=0.5mm of u\v] {\lab};
  }
  \node[factor] (f1) at (0.6,0)     {}; \draw[flink] (a1) -- (f1) -- (a2);
  \node[factor] (f2) at (1.8,0)     {}; \draw[flink] (a2) -- (f2) -- (a3);
  \node[factor] (f3) at (1.2,-0.7)  {}; \draw[flink] (a2) -- (f3) -- (a4);
  \node[factor] (f4) at (2.7,-0.7)  {}; \draw[flink] (a3) -- (f4) -- (a5);
  \node[prlab,above=0.4mm of f1] {$\psi_1$};
  \node[prlab,above=0.4mm of f2] {$\psi_2$};
  \node[prlab,left=0.4mm of f3]  {$\psi_3$};
  \node[prlab,right=0.4mm of f4] {$\psi_4$};
\end{tikzpicture}

\smallskip
\footnotesize
\setlength{\tabcolsep}{3.5pt}
\begin{tabular}{@{}lllcccc@{}}
\toprule
Factor & Scope & Relation & $00$ & $01$ & $10$ & $11$ \\
\midrule
$\phi_{1,\dots,5}$ & $a_i$ & prior $\frac12$ & \multicolumn{4}{c}{$a_i{=}0$: $\frac12$,\ \ $a_i{=}1$: $\frac12$} \\
\midrule
$\psi_1$ & $(a_1,a_2)$ & \ent, $p{=}.90$ & $1$ & $1$ & $\mathbf{.10}$ & $1$ \\
$\psi_2$ & $(a_2,a_3)$ & \ent, $p{=}.85$ & $1$ & $1$ & $\mathbf{.15}$ & $1$ \\
$\psi_3$ & $(a_2,a_4)$ & \con, $p{=}.90$ & $1$ & $1$ & $1$ & $\mathbf{.10}$ \\
$\psi_4$ & $(a_3,a_5)$ & \con, $p{=}.85$ & $1$ & $1$ & $1$ & $\mathbf{.15}$ \\
\bottomrule
\end{tabular}
\caption{The coherence MRF of Response~A. Circles are claims, grey squares the prior factors
  $\phi_i(x_i)$, black squares the relation factors. Each $\psi_r$ (Eq.~\ref{eq:psi}) is listed
  over the joint values $(x_s,x_t)$ of its two claims; bold cells are its violated set $V(\tau)$,
  down-weighted to $1-p_r$.}
\label{fig:mrf-a}
\end{figure}


Read directly, Eq.~\ref{eq:mrf} keeps each world's prior weight and multiplies it by $1-p_r$ for
every relation the world violates. A world that violates no relation keeps its prior weight, one
that violates an entailment of weight $0.9$ keeps a tenth of it, and with no relations $\PR=\pi$.

\paragraph{The check reading.} Picture a candidate world drawn at random from the prior. Each
relation acts as an independent \emph{check} on it. If the world breaks the relation, for example
by making the premise of an entailment true and its conclusion false, the check raises an alarm
with probability $p_r$; if the world respects the relation, it never does. A world therefore passes
the check of $r$ with probability $\psi_r(x_{e_r})$, and all checks with probability
$\prod_r\psi_r(x_{e_r})$. By Bayes' rule, the distribution of the surviving worlds is the prior
times this product, renormalized, which is exactly $\PR$. The normalizer $Z_\R$ is the probability
that a world drawn from the prior passes all checks.

For example, take two claims with flat priors, so each of the four worlds has probability $\frac14$,
and one entailment $a_1\Rightarrow a_2$ with $p_r=0.9$. Only the world $(1,0)$ breaks it, and it
survives the check with probability $0.1$. Hence $Z_\R=\frac34+\frac14\cdot0.1=0.775$, the world
$(1,0)$ falls to $0.025/0.775\approx0.03$, and each of the other three rises to
$0.25/0.775\approx0.32$. As $p_r\to1$ the check never misses, $(1,0)$ is eliminated, and the
relation becomes a hard constraint.

Figure~\ref{fig:mrf-a} shows the MRF of Response~A: five claims with flat priors, two entailment
factors that penalize a premise holding while its conclusion fails, and two contradiction factors
that penalize worlds in which a true claim and the falsehood it refutes both hold. A world drawn
from the priors passes all four checks with probability $Z_A=\vyAZ$. The committed world violates
none of them, and its probability rises from $\pi(\xs)=2^{-5}\approx\vyAPiX$ to $P_A(\xs)=\vyAPX$.
Under B's relations the same world falls to $P_B(\xs)=\vyBPX$, because it violates both of B's
conflicts.

\subsection{The Order and the Score (\S\ref{sec:score})}

\paragraph{The order.} Given the committed world, the order is deliberately minimal. It ranks only
edits that add or strengthen relations $\xs$ satisfies or remove or weaken relations $\xs$
violates, and it leaves edits that improve one part of an argument and degrade another unordered.
Arbitrary rewiring is covered whenever each step improves, and retyping a conflict as a support is
an improvement. The order is canonical relative to $\xs$, not independent of it: under verdict
polarity, adding an entailment between two false claims is an improvement, because it holds in the
committed world.

\paragraph{The axioms.} A2 makes a score invariant to presentation, such as reversing a symmetric
relation or swapping \sense{Precedence} for \sense{Succession}. It has a consequence worth stating:
relations on one pair that contradict each other, such as an exclusion and an equivalence of equal
weight, multiply to a constant and leave $\PR$ unchanged, so a score satisfying A2 rates them as
uninformative rather than as incoherent; we list this as a limitation.

\paragraph{The range of the score.} $\LCS$ lies in $(0,1)$, with three reference points
(Figure~\ref{fig:lcs-range}). A \emph{completely incoherent} response, whose relations rule out the
committed world, has $\LCS\to0$: $\PR(\xs)\to0$, as when $\xs$ violates conflicts of weight
approaching $1$ that a single other world satisfies all of. A response whose relations neither
confirm nor disconfirm $\xs$, including one with no relations, has $\PR(\xs)=\pi(\xs)$ and
$\LCS=\frac12$. A \emph{completely coherent} response, whose relations force the committed world,
has $\LCS\to1$, which requires relations that rule out every other world. Entailments, equivalences
and contradictions never rule out the world in which every claim is false, so forcing $\xs$ also
needs co-necessity or exclusion relations. Without them, and when $\xs$ is not that world, $n$
claims with flat priors stay below $\LCS=1-2^{-n}$, which is $31/32\approx0.97$ for five.

% The range of the LCS: LCS = sigma(W), W = logit P_R(x*) - logit pi(x*), the weight of
% evidence the relations give the committed world. Anchors: 0 (x* ruled out), 1/2
% (uninformative), 1 (x* forced). Voyager points from numbers.tex.
\begin{figure}[t]
\centering
\begin{tikzpicture}[x=5.9mm,y=36mm,font=\scriptsize]
  % axes
  \draw[-{Stealth[length=1.5mm]},black!70] (-5.6,0) -- (5.9,0) node[above left=0.3mm] {$W$};
  \draw[-{Stealth[length=1.5mm]},black!70] (-5.6,0) -- (-5.6,1.08) node[above] {$\LCS$};
  \foreach \y/\l in {0/0, 0.5/\frac12, 1/1}
    \draw[black!50] (-5.65,\y) -- (-5.55,\y) node[left=1.2mm,text=black] {$\l$};
  \foreach \x in {-4,-2,0,2,4}
    \draw[black!50] (\x,-0.02) -- (\x,0.02) node[below=1.5mm,text=black] {$\x$};
  \draw[dashed,black!40] (-5.6,1) -- (5.6,1);
  \draw[dotted,black!55] (-5.6,0.5) -- (0,0.5) -- (0,0);
  % the curve
  \draw[thick,blue!60!black,domain=-5.6:5.6,samples=100,smooth]
       plot (\x,{1/(1+exp(-\x))});
  % anchors
  \node[anchor=west,align=left,text=red!60!black] at (-5.4,0.36)
       {$\to0$: relations\\rule out $\xs$\\($\PR(\xs)\to0$)};
  \draw[red!60!black,-{Stealth[length=1.2mm]}] (-5.1,0.24) -- (-5.45,0.02);
  \node[anchor=east,align=right] at (-0.6,0.74) {$\frac12$: uninformative\\($\PR(\xs)=\pi(\xs)$)};
  \draw[black!70,-{Stealth[length=1.2mm]}] (-0.55,0.68) -- (-0.06,0.515);
  \node[anchor=east,align=right,text=blue!55!black] at (5.5,0.58)
       {$\to1$: relations force $\xs$\\($\PR(\xs)\to1$)};
  \draw[blue!55!black,-{Stealth[length=1.2mm]}] (5.1,0.66) -- (5.45,0.975);
  \fill[black] (0,0.5) circle (1pt);
  % Voyager responses
  \fill[blue!60!black] (\vyAW,\vyALCS) circle (1.4pt) node[above left=0.2mm] {A};
  \fill[red!60!black] (\vyBW,\vyBLCS) circle (1.4pt) node[below right=0.2mm] {B};
\end{tikzpicture}
\caption{The range of the score. $\LCS=\sigma(W)$, where
$W=\operatorname{logit}\PR(\xs)-\operatorname{logit}\pi(\xs)$, the argument of $\sigma$ in
Eq.~\ref{eq:lcs}, is the weight of evidence the relations give the committed world. $\LCS$ tends to $0$ when the relations rule $\xs$ out, equals
$\frac12$ when they leave its probability at the prior, and tends to $1$ when they force it.
Responses A and B of Figure~\ref{fig:voyager} score $\vyALCS$ and $\vyBLCS$.}
\label{fig:lcs-range}
\end{figure}


\paragraph{Two measures of confirmation.} The weight of evidence is also known as the
log-likelihood-ratio measure of confirmation \citep{fitelson1999}. In the check model, $\Lambda$ is
the log-ratio measure $\log P(\xs\mid E)/P(\xs)$ \citep{milne1996}, a different measure in general
\citep{fitelson1999}. Since $\PR(\xs)=\pi(\xs)e^{\Lambda}$, however, $\LCS$ is a strictly
increasing function of $\Lambda$ for a fixed claim set and prior (App.~\ref{app:lcs-lambda}). The
two differ in one respect. $\LCS$ depends on $\pi(\xs)$ and so on the whole claim set: adding a
claim that no relation touches leaves $\Lambda$ unchanged but moves $\LCS$ slightly towards
$\sigma(\Lambda)$. The additivity and bounds of Prop.~\ref{prop:props} are therefore properties of
$\Lambda$, not of $\LCS$. We nonetheless report $\LCS$: it is bounded, it reads as a probability,
with $\frac12$ as the uninformative point, and it stays informative when a pipeline's posteriors
are confident, where $\Lambda$ tends to $0$ for every graph (App.~\ref{app:twostage}).

\subsection{Why This Score (\S\ref{sec:why})}

\paragraph{Row-stochastic tables.} A natural alternative to Eq.~\ref{eq:psi} uses tables that are
row-stochastic in the source, so that a false premise leaves its conclusion at the prior
(App.~\ref{app:earlier}). Such tables fail Theorem~\ref{thm:necessity}'s test. At $\xs=\one$ they
are sign-definite with the intended sign only for $p\ge\frac12$. Below it the committed cell of an
entailment is the table's minimum, so a weak entailment lowers the score and a weak contradiction
raises it. When the source is false at $\xs$, the committed cell holds the prior's value
$\frac12$, which is neither the maximum nor the minimum, so they fail at every weight except
$p=\frac12$, where the table is constant and has no effect. For weight changes,
App.~\ref{app:proof-why} gives the analogue of Theorem~\ref{thm:necessity} as a sufficient
condition.

\paragraph{The readout world.} Proposition~\ref{prop:unique} does not say that $\xs$ is the right
world; choosing it is the factuality pipeline's job. It says that the order and the readout must
use the same world. It also does not exclude readouts that combine several worlds; of the five
App.~\ref{app:earlier} examines, four are not monotone and the fifth violates A2.

\paragraph{What is not forced.} These results fix the class of tables and the readout world; they
do not make $\Lambda$ the unique score satisfying A1--A4. Any score that is strictly increasing in
$\Lambda$ for a fixed prior and takes a prior-independent value at $\Lambda=0$ satisfies them too,
which is how $\LCS$ does. And any reweighting $p\mapsto1-(1-p)^\beta$ with $\beta>0$ gives another
monotone score of the same form, where different $\beta$ can rank incomparable graphs differently.
We take $\beta=1$ because it makes $p_r$ the probability that a check catches a violation, so the
check reading holds with the weights as estimated.

\paragraph{Why not count relations?} Monotonicity alone admits trivial scores: the signed weight sum
$\sum_{\Rsup}p_r-\sum_{\Rcon}p_r$ is monotone by construction. But a count credits every relation in
full, whether or not it rules out anything plausible, so it depends on how an argument is presented
rather than on what it rules out. Two entailments $a\Rightarrow b$ and $b\Rightarrow a$ of weight
$p$ multiply to exactly the table of one equivalence of weight $p$, so they induce the same $\PR$
and the same $\Lambda$, but they count $2p$ against $p$, which violates A2. Repeating a support
likewise adds its weight each time, where $\Lambda$ saturates (Prop.~\ref{prop:props}(vii)).
Averaging removes the dependence on the number of relations but breaks monotonicity, since a weak
support lowers an average.

\paragraph{Properties.} Additivity makes $\Lambda$ an amount of evidence, in \emph{nats}
(natural-log units: a change of $\Delta$ nats multiplies $\PR(\xs)$ by $e^{\Delta}$), that adds
across independent parts of an argument; $\LCS$ is not additive. A support earns credit only for
the probability it removes from worlds other than $\xs$, so once the relations force $\xs$, further
support earns almost nothing, and repeated copies of one relation converge rather than accumulate
(App.~\ref{app:proof-props}). The leave-one-out contributions $\Delta_r$ are not additive and need
not sum to $\Lambda$. Weighted mini-bucket elimination \citep{wmb}, for graphs beyond exact
inference, preserves the guarantees only up to the approximation.

\subsection{Logical Coherence and Other Notions of Coherence}
\label{sec:notions}

The definitions fix what the score responds to, and they separate logical coherence from the other
properties called coherence on five points. \citet{anonlocobench} tabulate the comparison; here we
state each point as a property of $\LCS$.
\begin{enumerate}[leftmargin=1.4em,itemsep=1pt,topsep=2pt]
\item \textbf{The text enters only through its relation graph.} $\LCS$ is a function of the
  relation graph, the priors and the committed world, so two texts with the same claims, verdicts
  and relations receive the same score, whatever their order and wording. Discourse coherence is a
  property of a text's organization: cohesion and local coherence track lexical and entity links
  between neighbouring sentences \citep{hallidayhasan1976,grosz1995}, and the standard test of a
  coherence model is that a shuffled text scores lower \citep{barzilay}.
\item \textbf{Only constraints on truth values count.} Relational theories of discourse coherence
  count any relation a reader can infer \citep{hobbs1979,mann}. The compile map sends
  \sense{Contrast}, \sense{Concession} and the temporal senses to \emph{none}, whose factor is
  constant. Conversely, a \sense{Refutation} between two true claims is well formed as discourse
  and a violated contradiction at $\xs$.
\item \textbf{Signs come from the committed world.} Discourse coherence, consistency and the
  coherence measures of Bayesian epistemology \citep{shogenji1999,olsson2002} do not depend on
  which claims are true. A relation's sign is fixed by $\xs$, so deriving a falsehood from a truth
  is a conflict, and refuting it is a support.
\item \textbf{Supports count, not only conflicts.} Consistency is the absence of contradiction
  \citep{summac,selfcheckgpt,mundler2024}, so it cannot rise when an inference is added. By
  Theorem~\ref{thm:mono} every added support of positive weight raises $\LCS$, and each further
  copy of the same support earns less (Prop.~\ref{prop:props}(vii)).
\item \textbf{The claims are held fixed, and their truth is not scored.} Probabilistic coherence
  measures compare different sets of propositions, and holistic ratings \citep{summeval,geval}
  compare whole texts that differ in content. The score's guarantee holds over one claim set and
  committed world, and a false claim affects it only through the relations it takes part in.
\end{enumerate}
Each point is tested by a ladder type of LoCoBench (\S\ref{sec-locobench}): the first by
\cp{order} and \cp{control}, the second by \cp{conflict}, the third by \cp{sign}, the fourth by
\cp{chain}, and the fifth by every ladder. These are choices of construct, not claims that
organization or style do not matter: logical coherence is meant to be reported beside discourse
coherence and factuality.

\paragraph{Relation to argumentation.} Gradual semantics for weighted bipolar argumentation
\citep{cayrol2005,amgoud2016,amgoud2018,baroni2019} also grade under axioms such as neutrality and
monotony, and probabilistic argumentation \citep{thimm2012,hunter2013} constrains a distribution by
attacks. Our setting differs in three ways. Relations are truth-functional couplings between claims
rather than attacks between arguments. A relation's sign is fixed by the factuality verdicts rather
than by its type. And monotonicity is asked of a readout of one probability model, so what must be
justified is the choice of potentials and of the world the readout uses.
